\chapter{Limitations}
\label{chap:limitations}

\section{Scope and Version}
\label{sec:limitations-scope}
\label{sec:limitations-version}

The current results describe the ordinary-list and finger-tree minimal cores
at source \texttt{4a83fa6}, with Python sources unchanged from calibration.
The earlier rollback-enabled experiment remains separate. The current matched
comparison measures those two complete sorting calls; it does not isolate
the causal effect of removing rollback. Update failure may leave a partial
state, which must be discarded. Successful-state invariants are not an
exception-recovery guarantee.

\section{Mathematical Model and Backend Fidelity}
\label{sec:limitations-backend}

The physical and locality arguments have passed this project's source-level
manual review. They have not been machine formalized or independently certified
by a separate external reviewer. The account makes endpoint and first-point
clarifications explicit, without claiming a unique verbatim interpretation
of every omitted paper detail. The finite-parent locality case is a separate
project argument; it is not mislabeled as the published latest-child lemma.

The tree's linear total bound uses unit-cost comparisons and machine-word
operations with expected dictionary lookup and amortized growth. Arbitrary
Python comparison methods, large-number bit costs, interpreter scheduling and
wall-clock behavior are outside it. Complete state audits and replay are also
outside the minimal complexity claim. A source change can invalidate the
review and its binding even if a few tests still pass.

Ordinary lists do not inherit that bound: insertion copying alone is quadratic
on increasing inputs, and their membership, ancestry and ownership work differs.
The current comparison therefore evaluates whole implementations. It cannot
attribute every difference solely to the abstract storage structure, nor use a
runtime curve as the proof of either complexity statement.

\section{Input Coverage and Scale}
\label{sec:limitations-sizes}
\label{sec:limitations-generality}

The common set contains fourteen controlled inputs at each of ten lengths
32/33 through 512/513. Its nesting strata have fixed quotas, and the inputs are
constructed rather than sampled from a population of Jordan sequences. Mirror
pairs and related constructors are dependent. Consequently, the overall ratio
is a descriptive median for this fixed set, not a population performance estimate.

All 3,788 recorded common-set splits per backend have an empty side, so that
set does not cover the interior split path. The independent targeted set
contains twelve constructions at 128/129, with nonempty smaller sides 3, 7 and
19. Its source size is always 61, its owner is finite, and its later direct
updates occur only on the retained side. It does not establish broad coverage
of all split ratios, black-height differences or both-output update histories.
Earlier operation tests and the 18-input complete-sorting campaign cover other
transitions (\cref{sec:implementation-coverage}), not additional performance
samples or large-scale mixed-update coverage on the measured version.

No measured input favors the tree. The narrowing common-set gap at larger
sizes does not guarantee a crossover beyond 513, or the same behavior for
another generator. Theoretical order bounds permit constants that dominate
throughout a finite measurement range.

\section{Timing Scopes and Environment}
\label{sec:limitations-timing-scopes}
\label{sec:limitations-execution}

The current interval includes fresh initialization, the shared algorithm,
output recovery and batching overhead. Both backends receive the same scope,
but their local checks, field-access wrappers and allocations differ. Their
individual contributions were not profiled or experimentally isolated.
Expected-output comparison, pre-batch GC collection and evidence persistence
are outside timing, while GC remains enabled during sorting.

The formal campaigns record one Apple M4 environment with CPython 3.12.14 and
battery power. Warm-ups and alternating first positions reduce selected startup
and order effects, but do not establish stable frequency, temperature or
absence of all unrelated operating-system activity. The campaigns' separate
identities and timestamps are retained. There is no multi-machine replication
of the current comparison.

Historical paper/reference ratios have a further scope mismatch: the reference
includes certification and its oracle-backed pipeline, while the old paper
call is pre-certified minimal sorting. Those ratios are not equivalent
end-to-end speedups. Historical CPython 3.12.4 observations, the development
pilot and current formal values are not pooled.

\section{Statistical and Exploratory Reach}
\label{sec:limitations-statistics}
\label{sec:limitations-exploratory}

Twenty batches characterize repeat variability for each fixed input/backend.
Dividing by B gives a per-call duration, not B independent samples. The median
and R7 IQR describe those observations; they do not remove persistent bias or
provide an independence justification for confidence intervals over related
inputs. Cross-input quartiles and ranges are distributions of case summaries,
not uncertainty intervals for a universal effect.

Some common-set cells have relative IQR above 50\%. They remain included,
without a post hoc noise filter, retry or parameter change. One sample batch
falls below the calibration target, which applied to three-batch calibration
medians rather than every future sample. The exact cause of either observation
was not measured. The frozen B table was used throughout.

Old within-size counter correlations are descriptive, based on an unbalanced
mix of twelve cases per size. Checked counters and minimal times also come
from different observation policies. These correlations cannot causally
explain the current tree/list difference. In particular, a constant invariant
counter produces an undefined correlation rather than a zero effect.

\section{Output Provenance and Recognition}
\label{sec:limitations-provenance}

The current paper cores traverse their own maintained output list. External
certification and expected sorting are used only for qualification and output
checking. The distinct reference pipeline deliberately returns oracle-derived
output; that provenance must not be attributed to the paper core.

All formal inputs are pre-certified valid. The evidence does not establish
recognition accuracy, invalid-input failure reasons or recognition runtime.
The sorting proof likewise starts from its validity precondition. Error
correction and geometric applications remain outside this evaluation.

\section{Consequences for the Conclusions}
\label{sec:limitations-conclusions}

The thesis can report a source-reviewed structural bound in an explicit model,
finite qualification and output success, and the measured tree slowdown for
two frozen sets. It cannot infer an unconditional Python time bound, a practical
speed advantage, independent random sampling, a future crossover, a causal
cost decomposition or hardware-independent absolute runtime. These boundaries
preserve the distinction between the implemented algorithm, its mathematical
account and its measured behavior.
